\documentclass[11pt,reqno]{amsart}

\usepackage{amsmath,amssymb,amsthm}
\usepackage[margin=1.15in]{geometry}
\usepackage[colorlinks=true,linkcolor=blue,citecolor=blue]{hyperref}

\newtheorem{theorem}{Theorem}[section]
\newtheorem{proposition}[theorem]{Proposition}
\newtheorem{lemma}[theorem]{Lemma}
\newtheorem{corollary}[theorem]{Corollary}
\theoremstyle{definition}
\newtheorem{definition}[theorem]{Definition}
\newtheorem{question}[theorem]{Question}
\newtheorem{remark}[theorem]{Remark}

\newcommand{\HH}{\mathcal{H}}
\newcommand{\PP}{\mathcal{P}}
\newcommand{\N}{\mathbb{N}}
\newcommand{\R}{\mathbb{R}}
\newcommand{\C}{\mathbb{C}}
\newcommand{\U}{\mathcal{U}}
\newcommand{\FF}{\mathcal{F}}
\newcommand{\supp}{\operatorname{supp}}
\newcommand{\lk}{\ell^2(\kappa)}
\newcommand{\Blk}{\mathcal{B}(\ell^2(\kappa))}
\newcommand{\sess}{\sigma_{\mathrm{ess}}}

\title[Projection-lattice ultrafilters at a measurable cardinal]
{Projection-lattice ultrafilters\\ at a measurable cardinal}

\author{David Feldman}
\address{Department of Mathematics and Statistics, University of New Hampshire, Durham, NH}

\author{Alexander Wilce}
\address{Department of Mathematics, Susquehanna University, Selinsgrove, PA}

\date{\today}

\begin{document}

\begin{abstract}
Let $\kappa$ be a measurable cardinal, $\U$ a normal measure on $\kappa$, and
$(e_\alpha)_{\alpha<\kappa}$ an orthonormal basis of $\lk$. Because every vector of
$\lk$ has countable support and $\U$ is closed under diagonal intersections, every
bounded operator becomes \emph{exactly} diagonal upon compression to $\ell^2(S)$ for a
suitable $S\in\U$: Anderson paving with error zero, by a diagonal intersection. We
deduce an exact dichotomy --- every closed subspace, on a measure-one set of
coordinates, either contains the coordinate block or misses it entirely --- and
conclude that the block filter of a \emph{normal} measure is already a maximal filter
in the projection lattice (a \emph{pluf}), $\kappa$-complete and nonprincipal ---
normality, not mere $\kappa$-completeness, is what the diagonal-intersection argument
consumes, exactly as selectivity, not mere ultrafilterhood, is what maximality
consumes at $\aleph_0$. The associated state is
a singular, $<\kappa$-additive pure state. This yields proofs of the existence
directions of two theorems of Blecher and Weaver (singular countably additive,
respectively $<\kappa$-additive, pure states on $\Blk$ exist when $\kappa$ is Ulam
measurable, respectively measurable) that make no use of the Marcus--Spielman--Srivastava
solution of the Kadison--Singer problem. In the converse direction, Blecher--Weaver's
excision theory shows that the $1$-set of \emph{any} countably additive pure state on
$\mathcal B(H)$ is a pluf; hence quantum-filter maximality and lattice maximality, which
diverge for the singular pure states available below the first Ulam measurable cardinal,
coincide above it.
\end{abstract}

\maketitle

\section{Introduction}\label{sec:intro}

This note is the large-cardinal chapter of a study of maximal filters in the projection
lattice of a Hilbert space, carried out in the companion papers \cite{FWI,FWII}. There,
on a separable space, the maximality of the filter generated by the coordinate
subspaces of an ultrafilter turned on delicate matters: it holds exactly for selective
ultrafilters \cite{FWII}, the uniqueness of the associated state is the Kadison--Singer
problem \cite{KadisonSinger}, resolved by Marcus, Spielman and Srivastava \cite{MSS},
and under CH maximal filters exist along which no ellipsoid has round slices
\cite{FWI}. On a Hilbert space of measurable dimension these questions change
character, by a mechanism with two ingredients: vectors of $\lk$ have countable
support, and normal measures absorb diagonal intersections. The rest of the paper
follows from the resulting Theorem~\ref{thm:F}.

Our reference point throughout is the paper of Blecher and Weaver
\cite{BlecherWeaver}, which characterizes the cardinals $\kappa$ such that $\Blk$
carries singular states with various additivity properties: singular countably additive
pure states exist iff $\kappa$ is Ulam measurable, and singular $<\kappa$-additive pure
states iff $\kappa$ is measurable. Their existence proofs route purity through a
generalization of Kadison--Singer to arbitrary $\kappa$ (\cite{BlecherWeaver}, Theorem
5.1), obtained by gluing Marcus--Spielman--Srivastava pavings over all countable subsets
of $\kappa$ by compactness. Section~\ref{sec:mssfree} shows that for the existence
statements this heavy input can be dispensed with: along a normal measure, purity is a
diagonal intersection. The full Blecher--Weaver Theorem 5.1 --- unique extension for
\emph{arbitrary} pure states on the diagonal, including those living on countable sets
of coordinates, where the problem is verbatim the classical one --- of course still
requires \cite{MSS}; the boundary between soft and hard runs through the completeness
of the ultrafilter, not through the cardinal.

Section~\ref{sec:J} runs the bridge in the other direction. On a separable space, the
$1$-set $\FF_\varphi=\{p:\varphi(p)=1\}$ of a singular pure state is a maximal
\emph{quantum filter} in the sense of Farah--Weaver \cite{FarahWofsey} but never a
lattice filter: it is not closed under meets, which is precisely why the
non-diagonalizable pure states of Akemann--Weaver and Koszmider produce no maximal
lattice filters \cite{FWI}. Blecher--Weaver's regularity and excision theory
(\cite{BlecherWeaver}, \S4) shows that countable additivity repairs this completely:
the $1$-set of a countably additive pure state is a pluf (Theorem~\ref{thm:J}). Since
singular countably additive pure states exist exactly at Ulam measurable dimension, the
divergence of quantum-filter maximality from lattice maximality is a phenomenon of
small dimension.

\subsection*{Conventions}
$\kappa$ is an uncountable cardinal; a \emph{normal measure} on $\kappa$ is a
$\kappa$-complete nonprincipal ultrafilter $\U$ closed under diagonal intersections:
if $X_\alpha\in\U$ for $\alpha<\kappa$ then
$\triangle_\alpha X_\alpha=\{\beta:\beta\in X_\alpha\text{ for all }\alpha<\beta\}
\in\U$. Every measurable cardinal carries one. We fix an orthonormal basis
$(e_\alpha)_{\alpha<\kappa}$ of $\HH=\lk$; for $S\subseteq\kappa$,
$\HH_S=\overline{\operatorname{span}}\{e_\alpha:\alpha\in S\}$. A \emph{pluf} is a
maximal filter in the lattice $\PP(\HH)$ of closed subspaces; a pluf is
\emph{principal} if it consists of the closed subspaces containing a fixed line, and by
\cite{FWI}, Lemma 2.2, a pluf containing any finite-dimensional subspace is principal.
A filter is maximal iff every closed subspace either belongs to it or meets some member
trivially (\cite{FWI}, Lemma 2.1). A state on $\mathcal B(H)$ is \emph{singular} if it
annihilates the compact operators; $<\kappa$-additivity of states is as in
\cite{BlecherWeaver}.

\section{Exact paving by diagonal intersection}\label{sec:paving}

\begin{theorem}\label{thm:F}
Let $\U$ be a normal measure on $\kappa$ and $T\in\Blk$. Then there is $S\in\U$ such
that $P_STP_S$ is diagonal: $\langle Te_\alpha,e_\beta\rangle=0$ for all distinct
$\alpha,\beta\in S$.
\end{theorem}

\begin{proof}
Each vector of $\lk$ has countable support, so
$c_\alpha:=\supp(Te_\alpha)\cup\supp(T^*e_\alpha)$ is countable and
$\kappa\setminus c_\alpha\in\U$ by $\kappa$-completeness. Let
$S=\triangle_{\alpha<\kappa}(\kappa\setminus c_\alpha)\in\U$. For $\alpha<\beta$ in
$S$ we have $\beta\notin c_\alpha$, whence $\langle Te_\alpha,e_\beta\rangle=0$ and
$\langle Te_\beta,e_\alpha\rangle=\langle e_\beta,T^*e_\alpha\rangle
=\overline{\langle T^*e_\alpha,e_\beta\rangle}=0$.
\end{proof}

This is Anderson's paving \cite{Anderson79} with error zero: not small norm on
finitely many diagonal pieces, but vanishing off-diagonal on a single measure-one
piece. Its spectral consequence:

\begin{corollary}\label{cor:flat}
Let $T\in\Blk$, $\epsilon>0$, and
$L=\lim_{\U}\langle Te_\alpha,e_\alpha\rangle$ (the $\U$-limit exists since $T$ is
bounded). Then there is $S\in\U$ with
\[
\bigl|\langle Tx,x\rangle-L\|x\|^2\bigr|\le\epsilon\|x\|^2
\qquad\text{for all }x\in\HH_S .
\]
For self-adjoint $T$ this says $T_{\HH_S}$ is diagonal with all eigenvalues in
$[L-\epsilon,L+\epsilon]$.
\end{corollary}

\begin{proof}
Intersect the $S$ of Theorem~\ref{thm:F} with
$\{\alpha:|\langle Te_\alpha,e_\alpha\rangle-L|\le\epsilon\}\in\U$. For
$x=\sum_{\alpha\in S}c_\alpha e_\alpha$ the off-diagonal terms of
$\langle Tx,x\rangle$ vanish, leaving
$\sum_\alpha|c_\alpha|^2\langle Te_\alpha,e_\alpha\rangle$, which lies within
$\epsilon\sum_\alpha|c_\alpha|^2$ of $L\|x\|^2$. Self-adjointness is nowhere used:
the quadratic form sees only the symmetric part of $T$.
\end{proof}

\section{The subspace dichotomy and maximality}\label{sec:dichotomy}

\begin{theorem}\label{thm:G}
Let $\U$ be a normal measure on $\kappa$. For every closed subspace $W\subseteq\lk$
there is $S\in\U$ with either $\HH_S\subseteq W$ or $W\cap\HH_S=0$.
\end{theorem}

\begin{proof}
Apply Theorem~\ref{thm:F} to $T=P_{W^\perp}$, obtaining $S_0\in\U$ with
$\langle P_{W^\perp}e_\alpha,e_\beta\rangle=0$ for distinct $\alpha,\beta\in S_0$.
Write $g_\alpha=P_{W^\perp}e_\alpha$; then for distinct $\alpha,\beta\in S_0$,
\[
\langle g_\alpha,g_\beta\rangle
=\langle P_{W^\perp}e_\alpha,P_{W^\perp}e_\beta\rangle
=\langle P_{W^\perp}e_\alpha,e_\beta\rangle=0:
\]
the family $(g_\alpha)_{\alpha\in S_0}$ is \emph{orthogonal}. Hence for
$x=\sum_{\alpha\in S_0}c_\alpha e_\alpha\in\HH_{S_0}$,
\[
\|P_{W^\perp}x\|^2=\sum_{\alpha\in S_0}|c_\alpha|^2\|g_\alpha\|^2,
\]
with no cross terms, so $x\in W$ iff $c_\alpha=0$ whenever $g_\alpha\neq0$. Writing
$S_1=\{\alpha\in S_0:g_\alpha=0\}=\{\alpha\in S_0:e_\alpha\in W\}$, this says
$W\cap\HH_{S_0}=\HH_{S_1}$ exactly. Now $\U$ decides: if $S_1\in\U$, take $S=S_1$ and
$\HH_S\subseteq W$; otherwise $S=S_0\setminus S_1\in\U$ and $W\cap\HH_S=0$.
\end{proof}

\begin{remark}
On a measure-one set, incidence and angle coincide: $W$ and the coordinate blocks are
in exact general position. This is the geometric content of the theorem, and the reason
the pathologies of \cite{FWI,FWII} --- all of which stem from subspaces meeting a block
trivially while lying at angle zero to it --- vanish along $\U$. The contrast with the
separable dichotomy of \cite{FWII}, Corollary 2.2 is sharp: there the surviving
alternative is \emph{finite codimension in} $\HH_S$; here it is containment of $\HH_S$
on the nose, and no auxiliary finite-codimension subspaces need be adjoined to reach
maximality.
\end{remark}

\begin{corollary}\label{cor:H}
$\Phi(\U):=\{M\in\PP(\lk):\HH_S\subseteq M\text{ for some }S\in\U\}$ is a nonprincipal,
$\kappa$-complete pluf.
\end{corollary}

\begin{proof}
$\Phi(\U)$ is a proper filter: $\HH_S\cap\HH_{S'}=\HH_{S\cap S'}$ and $\U$ is a filter
of nonempty sets; upward closure is built in; $\kappa$-completeness of $\U$ gives
closure under intersections of fewer than $\kappa$ members (the intersection of blocks
over sets $S_i\in\U$, $i<\lambda<\kappa$, contains
$\HH_{\bigcap S_i}$). Maximality: given closed $W$, Theorem~\ref{thm:G} either puts
$W\in\Phi(\U)$ or produces $\HH_S\in\Phi(\U)$ with $W\cap\HH_S=0$; apply \cite{FWI},
Lemma 2.1. Nonprincipality: were $\Phi(\U)$ principal at a line $L$, then $L\subseteq
\HH_S$ for all $S\in\U$; but $L$ has countable support $c$, and
$\kappa\setminus c\in\U$ gives $L\cap\HH_{\kappa\setminus c}=0$.
\end{proof}

\begin{corollary}\label{cor:I}
Define $\varphi_\U(T)=\lim_{\U}\langle Te_\alpha,e_\alpha\rangle$. Then
$S_{\Phi(\U)}=\{\varphi_\U\}$ --- the state-space face of \cite{FWI}, \S3 is a
singleton --- so $\varphi_\U$ is a pure state, and $\Phi(\U)$ has the round slice
property for every ellipsoid: for every positive invertible $T$ and $\epsilon>0$ some
$M\in\Phi(\U)$ has $r(E_T\cap M)<1+\epsilon$, with
$\lim_{\Phi(\U)}T=\varphi_\U(T)$.
\end{corollary}

\begin{proof}
Let $\varphi\in S_{\Phi(\U)}$ and $T$ self-adjoint. By Corollary~\ref{cor:flat} pick
$S\in\U$ with $(L-\epsilon)P_{\HH_S}\le P_{\HH_S}TP_{\HH_S}\le(L+\epsilon)P_{\HH_S}$
where $L=\varphi_\U(T)$; since $\varphi(P_{\HH_S})=1$,
$\varphi(T)=\varphi(P_{\HH_S}TP_{\HH_S})\in[L-\epsilon,L+\epsilon]$. As $\epsilon$ was
arbitrary, $\varphi(T)=L$ for every self-adjoint $T$: the face is the singleton
$\{\varphi_\U\}$, whose element, being an extreme point of the state space
(\cite{FWI}, Proposition 3.4), is pure. The round-slice statements are the case of
positive invertible $T$, via \cite{FWI}, Lemma 3.1 and Proposition 3.2.
\end{proof}

\begin{remark}[Scope]
The softness is local to the normal measure. An ultrafilter on $\kappa$ concentrating
on a countable subset reproduces the classical Kadison--Singer problem verbatim, and
unique extension for arbitrary pure states on the diagonal of $\Blk$ ---
Blecher--Weaver's Theorem 5.1 --- still requires \cite{MSS}. Likewise the general
pathology of \cite{FWI,FWII} persists at $\kappa$: the pluf space over $\lk$ is
non-compact, has no prime filters, and admits non-intimate subspaces (graphs of
unitaries between complementary blocks), so non-diagonalizable plufs are not excluded.
What Theorem~\ref{thm:G} \emph{does} exclude is any witness family for the reduction
of \cite{FWII}, Proposition 5.3 consisting of coordinate subspaces of any single
basis.
\end{remark}

\begin{remark}[Normality is the $\kappa$-analogue of selectivity]\label{rem:normality}
Corollary~\ref{cor:H} consumes normality, and the parallel with the separable theory
is exact. At $\aleph_0$, the block
filter of an ultrafilter (with the finite-codimension subspaces adjoined) is maximal
if and only if the ultrafilter is \emph{selective} (\cite{FWII}, Theorem 3.4) --- and
selective ultrafilters are exactly the Rudin--Keisler minimal ones. At $\kappa$, the
proof of Theorem~\ref{thm:F} uses the diagonal intersection, i.e.\ \emph{normality},
not mere $\kappa$-completeness --- and normal ultrafilters are exactly the
Rudin--Keisler minimal ones among $\kappa$-complete ultrafilters on $\kappa$, with
the Rowbottom partition property playing the role of the Ramsey property. So the
same combinatorial regularity (RK-minimality in the appropriate category) powers
maximality at both cardinals; the difference is that at $\kappa$ every measurable
cardinal supplies a normal measure outright, while at $\aleph_0$ selectivity is an
extra, independence-prone hypothesis. How much of the converse direction of
\cite{FWII}, Theorem 3.4 survives at $\kappa$ is settled in part by the next two
results. Every vector of $\lk$ --- hence every constraint vector of a witness
subspace --- lives on countably many coordinates, so the witness construction is
governed entirely by partitions of $\kappa$ into \emph{countable} pieces, which are
automatically $\U$-small by $\kappa$-completeness. The fibers of a regressive function
play no role; the pieces simply \emph{are} the supports.
\end{remark}

\begin{definition}
Say a $\kappa$-complete nonprincipal ultrafilter $\U$ on $\kappa$ is a
\emph{$\sigma$-Q-point} if every partition of $\kappa$ into countable nonempty
pieces admits a partial selector in $\U$ (a set meeting each piece at most once);
equivalently, if every countable-to-one function $g:\kappa\to\kappa$ is injective
on a set of $\U$. (For the equivalence: the fibers of a countable-to-one $g$ form
such a partition, and injectivity on $S$ is exactly the selector property; in the
other direction, send each point to the least element of its piece --- a
countable-to-one function whose injectivity sets are the selectors.) This is the
exact $\kappa$-analogue of the Q-point property at $\aleph_0$, with
``finite-to-one'' promoted to ``countable-to-one''.
\end{definition}

\begin{lemma}\label{lem:normal-sigmaQ}
Every normal measure is a $\sigma$-Q-point. Indeed, for a partition of $\kappa$
into countable pieces, the set $M$ of piece-minima itself lies in $\U$, and $M$
is a selector.
\end{lemma}

\begin{proof}
Suppose $M\notin\U$, so $X=\kappa\setminus M\in\U$. The map
$g(\alpha)=\min(\text{piece of }\alpha)$ satisfies $g(\alpha)<\alpha$ on $X$
(a point outside $M$ is not the minimum of its piece), so by normality
(Fodor) $g$ is constant on a set of $\U$, which is then contained in a single
countable piece --- impossible for a $\kappa$-complete nonprincipal ultrafilter.
So $M\in\U$; and $|M\cap\sigma|=1$ for every piece $\sigma$.
\end{proof}

\begin{proposition}\label{prop:kappa-witness}
Let $\U$ be $\kappa$-complete and nonprincipal. If $\U$ is \emph{not} a
$\sigma$-Q-point, then $\Phi(\U)$ is not maximal.
\end{proposition}

\begin{proof}
Fix a partition $\kappa=\bigsqcup_i\sigma_i$ into countable pieces with no
partial selector in $\U$. For each $i$ enumerate $\sigma_i=\{\alpha_{i,n}:n<
|\sigma_i|\}$ and set $f_i=\sum_n2^{-n-1}e_{\alpha_{i,n}}\in\lk$: square-summable
because the piece is countable, with full support $\sigma_i$. Let $W=\{f_i:i\}^\perp$.

Since the pieces are pairwise disjoint and cover $\kappa$, every $x\in\HH_S$
decomposes as $x=\sum_ix^{(i)}$ with $x^{(i)}\in\ell^2(\sigma_i\cap S)$, and
$x\in W$ iff each $\langle x^{(i)},f_i\rangle=0$. Three consequences. First,
$W\notin\Phi(\U)$: membership requires $\HH_S\subseteq W$ for some $S\in\U$,
i.e.\ $S$ disjoint from every support --- but the supports cover $\kappa$.
Second, no member of $\Phi(\U)$ blocks $W$: a member contains some $\HH_S$,
$S\in\U$, and $W\cap\HH_S=0$ holds iff for every $i$ the only vector of
$\ell^2(\sigma_i\cap S)$ orthogonal to $f_i$ is zero, iff $|\sigma_i\cap S|\le1$
for every $i$ (the coefficients of $f_i$ are nonzero throughout $\sigma_i$) ---
that is, iff $S$ is a partial selector, excluded by hypothesis. So
$W\cap\HH_S\neq0$ for every $S\in\U$, and third, $W$ is therefore
\emph{addable}: it meets every member of $\Phi(\U)$. Addable and unadded:
by the maximality criterion (\cite{FWI}, Lemma 2.1), $\Phi(\U)$ is not maximal.
\end{proof}

\begin{corollary}\label{cor:necessity}
$\Phi(\U)$ maximal $\Rightarrow$ $\U$ is a $\sigma$-Q-point; and normal
$\Rightarrow$ $\sigma$-Q-point (Lemma~\ref{lem:normal-sigmaQ}). The
characterization question is thus pinched between the two: whether the
$\sigma$-Q-point property is also \emph{sufficient} for maximality, and whether
it is strictly weaker than isomorphism to a normal measure, are
Question~\ref{q:kappaC}.
\end{corollary}

\begin{remark}[Formal verification]\label{rem:lean3}
Every lattice-level assertion of this paper has been formalized and
machine-checked in Lean~4 with Mathlib (toolchain v4.28.0):
Theorem~\ref{thm:F} (exact paving); Corollary~\ref{cor:flat} in the
quadratic-form statement given above, together with the existence of ultrafilter
limits of bounded functions;
Theorem~\ref{thm:G} (the exact dichotomy); Corollary~\ref{cor:H} in full
filter language (the block filter is proper, upward closed, closed under
intersections --- completeness being certified generically in the index type, so that
$\sigma$- and ${<}\kappa$-completeness are instances --- nonprincipal, and decides
every closed subspace);
Lemma~\ref{lem:normal-sigmaQ} in the abstract form (Fodor property plus the
absence of countable measure-one sets imply the $\sigma$-Q-point property,
the piece-minima forming a measure-one selector); and
Proposition~\ref{prop:kappa-witness} in both halves --- the combinatorial
($\sigma$-Q equivalences and the transversal lemma) and the Hilbert-space
(the witness subspace of a countable-piece partition, unadded yet meeting
every member of the block filter). The mechanized development also certifies
that a block meets the witness trivially exactly when its index set is a partial
selector; the hypothesis that the pieces cover $\kappa$ is necessary there, and a
counterexample without it is certified alongside. The artifacts compile with no
unproved obligations, and every statement's axiom footprint is exactly
\texttt{propext}, \texttt{Classical.choice}, \texttt{Quot.sound}. No
large-cardinal machinery is axiomatized: closure under diagonal
intersections, countable completeness, and the smallness of countable sets
enter only as hypotheses, so the mechanized statements are theorems of ZFC
about arbitrary ultrafilters on a (for the combinatorics, well-ordered)
index type. Sections~\ref{sec:mssfree}--\ref{sec:J} are likewise machine-checked, with
the excision and regularity theory of \cite{BlecherWeaver} quarantined as a
named package of hypotheses --- $\sigma$-filtration of the $1$-set and the
existence of excising projections --- and with no appeal, direct or indirect,
to \cite{MSS}: the ultrafilter-limit state is verified to be a state, to
annihilate the compact operators, to be countably additive (and additive in
the generic-index form that specializes to $<\kappa$-additivity), and to be
the unique element of the face of $\Phi(\U)$, whence pure; the $1$-set of a
state satisfying the excision package is verified to be a pluf, nonprincipal
when the state annihilates the finite-rank projections; and the $1$-set of
the limit state is verified to be exactly $\Phi(\U)$. The formalization is
carried out over the reals throughout; no argument of these sections required
complex scalars. The verification artifacts will be archived with the final
version of this paper.
\end{remark}

\section{Existence of singular additive pure states, without paving}\label{sec:mssfree}

\begin{theorem}\label{thm:mssfree}
Let $\kappa_0$ be measurable with normal measure $\U$, and let $\kappa\ge\kappa_0$.
Regard $\ell^2(\kappa_0)\subseteq\lk$ as a coordinate block. Then
\[
\varphi(T)\;=\;\lim_{\alpha\to\U}\langle Te_\alpha,e_\alpha\rangle
\qquad(T\in\Blk)
\]
is a singular, countably additive pure state on $\Blk$; for $\kappa=\kappa_0$ it is
$<\kappa_0$-additive. In particular the existence directions of Theorems 2.4(ii) and
2.4(iv) of \cite{BlecherWeaver} hold with no appeal to the Kadison--Singer theorem of
\cite{MSS}.
\end{theorem}

\begin{proof}
$\varphi$ is a state (a $\U$-limit of vector states). Diagonal intersections over
$\alpha<\kappa_0$ of the countable supports $c_\alpha\subseteq\kappa$ go through
verbatim ($c_\alpha\cap\kappa_0$ is countable, hence $\kappa_0$-small), so
Theorems~\ref{thm:F} and \ref{thm:G} hold for blocks indexed by $S\in\U\subseteq
\mathcal P(\kappa_0)$ inside $\PP(\lk)$, and Corollaries~\ref{cor:H} and \ref{cor:I}
give: $\Phi(\U)$ is a pluf in $\PP(\lk)$ with $S_{\Phi(\U)}=\{\varphi\}$, so $\varphi$
is pure.

\emph{Singularity.} A rank-one projection $P_v$ has
$\langle P_ve_\alpha,e_\alpha\rangle=|\langle v,e_\alpha\rangle|^2$, nonzero for only
countably many $\alpha$; $\kappa_0$-completeness gives $\varphi(P_v)=0$, and
singularity follows since singular states are exactly those vanishing on all rank-one
projections.

\emph{Countable additivity.} Let $(q_n)$ be orthogonal projections. At most countably
many $\varphi(q_n)$ are nonzero and $\sum\varphi(q_n)\le1$; splitting off that part and
using finite additivity, it suffices to treat $\varphi(q_n)=0$ for all $n$. Given
$\epsilon>0$, the sets
$A_n=\{\alpha:\langle q_ne_\alpha,e_\alpha\rangle<\epsilon2^{-n}\}$ lie in $\U$
(their complements would give $\varphi(q_n)\ge\epsilon2^{-n}$), and
$\bigcap_nA_n\in\U$ by countable completeness; on it
$\langle(\sum q_n)e_\alpha,e_\alpha\rangle=\sum_n\langle
q_ne_\alpha,e_\alpha\rangle<2\epsilon$, so $\varphi(\sum q_n)\le2\epsilon$ for all
$\epsilon$, i.e.\ $=0$.

\emph{$<\kappa_0$-additivity when $\kappa=\kappa_0$.} By \cite{BlecherWeaver},
Proposition 2.2(iii) it suffices to check normality of $\varphi$ on atomic abelian
subalgebras generated by fewer than $\kappa_0$ orthogonal projections, which is the
transfinite monotone convergence Lemma 3.2 of \cite{BlecherWeaver} applied to the
$\kappa_0$-additive measure $\U$; alternatively, run the argument of the previous
paragraph with $(q_\beta)_{\beta<\lambda}$, $\lambda<\kappa_0$, replacing the sets
$A_n$ by
$A^\epsilon_\beta=\{\alpha:\langle q_\beta e_\alpha,e_\alpha\rangle<\epsilon_\beta\}$
for any summable $(\epsilon_\beta)$ supported on the countably many relevant $\beta$
--- for each fixed $\alpha$, $\langle q_\beta e_\alpha,e_\alpha\rangle\neq0$ for only
countably many $\beta$, and $\kappa_0$-completeness applies.
\end{proof}

\begin{remark}
For Ulam measurable $\kappa$ (i.e.\ $\kappa\ge$ some measurable $\kappa_0$) the theorem
produces the singular countably additive pure state on $\Blk$ required by
\cite{BlecherWeaver}, Theorem 2.4(ii); the reverse directions of the Blecher--Weaver
characterizations, being nonexistence statements, are untouched and remain as proved
there. What the diagonal intersection replaces is exactly the compactness gluing of
uncountably many Marcus--Spielman--Srivastava pavings in their Theorem 5.1 --- for
those pure states on the diagonal that come from $\kappa$-complete ultrafilters. The
hard/soft boundary in Kadison--Singer at large $\kappa$ thus runs through completeness,
not cardinality.
\end{remark}

\section{Countably additive pure states give plufs}\label{sec:J}

We now use Blecher--Weaver's own machinery to close the circle. Recall
(\cite{BlecherWeaver}, \S4): a state $\varphi$ on a von Neumann algebra is
\emph{regular} if countable joins of null projections are null; every countably
additive state is regular, and regular \emph{pure} states are countably additive
(\cite{BlecherWeaver}, Theorem 4.1, after \cite{Hamhalter}); for regular $\varphi$ the
$1$-set $\FF_\varphi=\{p:\varphi(p)=1\}$ is a $\sigma$-filter (\cite{BlecherWeaver},
Lemma 4.3); a projection $p\in\FF_\varphi$ \emph{excises} $\varphi$ for $x$ if
$pxp=\varphi(x)p$, and for countably additive pure $\varphi$ every $x$ admits an
excising projection in $\FF_\varphi$ (\cite{BlecherWeaver}, Lemma 4.5); finally, a pure
state is determined by its $1$-set (\cite{BlecherWeaver}, Lemma 4.2).

\begin{theorem}\label{thm:J}
Let $\varphi$ be a countably additive pure state on $\mathcal B(H)$, $H$ of arbitrary
dimension. Then $\FF_\varphi$ is a pluf (identifying projections with their ranges),
and $S_{\FF_\varphi}=\{\varphi\}$. If $\varphi$ is singular, $\FF_\varphi$ is
nonprincipal.
\end{theorem}

\begin{proof}
$\FF_\varphi$ is a filter: closed under (countable) meets by regularity
(\cite{BlecherWeaver}, Lemma 4.3), upward closed, and proper since
$\varphi(0)=0$. Maximality: let $M$ be a closed subspace with $P_M\notin\FF_\varphi$,
i.e.\ $c:=\varphi(P_M)<1$. By excision (\cite{BlecherWeaver}, Lemma 4.5) choose
$p\in\FF_\varphi$ with $pP_Mp=cp$. If $v$ were a unit vector of
$\operatorname{ran}p\cap M$, then
$1=\langle P_Mv,v\rangle=\langle pP_Mpv,v\rangle=c<1$, absurd; so
$\operatorname{ran}p\cap M=0$, and \cite{FWI}, Lemma 2.1 gives maximality. That
$S_{\FF_\varphi}=\{\varphi\}$ is \cite{BlecherWeaver}, Lemma 4.2. If $\varphi$ is
singular it vanishes on every rank-one projection, so $\FF_\varphi$ contains no line;
a principal pluf contains a line (\cite{FWI}, Lemma 2.2), so $\FF_\varphi$ is
nonprincipal.
\end{proof}

\begin{corollary}\label{cor:coincide}
For countably additive pure states on $\mathcal B(H)$, the maximal quantum filter
$\FF_\varphi$ of Farah--Weaver \cite{FarahWofsey} \emph{is} a maximal lattice filter:
quantum-filter maximality and lattice maximality coincide. By contrast, for singular
pure states on a separable $H$ --- which are never countably additive --- the $1$-set
is not even a filter (\cite{FarahWofsey}; \cite{BlecherWeaver}, remark before Lemma
4.3), which is why the non-diagonalizable pure states of \cite{AW,Koszmider} yield no
plufs \cite{FWI}. Since singular countably additive pure states exist exactly when
$\dim H$ is Ulam measurable (\cite{BlecherWeaver}, Theorem 2.4(ii)), the divergence of
the two maximality notions for singular states is a phenomenon of dimension below the
first Ulam measurable cardinal.
\end{corollary}

\begin{corollary}\label{cor:Fphi-is-Phi}
In the setting of Corollary~\ref{cor:I}, $\FF_{\varphi_\U}=\Phi(\U)$: the pluf is
exactly the $1$-set of its state.
\end{corollary}

\begin{proof}
$\Phi(\U)\subseteq\FF_{\varphi_\U}$ since $\varphi_\U(P_M)=1$ for $M\in\Phi(\U)$;
both are plufs (Corollary~\ref{cor:H}, Theorem~\ref{thm:J}); a maximal filter admits
no proper filter extension.
\end{proof}

\section{Questions}\label{sec:questions}

\begin{question}\label{q:anderson}
Anderson's conjecture at $\kappa$: is every countably additive pure state on $\Blk$ of
the form $\varphi_\U\circ$(basis change) for some $\kappa$-complete ultrafilter and
orthonormal basis --- equivalently (via Theorem~\ref{thm:J} and \cite{FWII}, Theorem
5.1), is every pluf of the form $\FF_\varphi$, $\varphi$ countably additive,
diagonalizable? Blecher--Weaver note the conjecture is expected to be consistent; a
counterexample would produce a non-diagonalizable pluf with singleton state face,
answering the quadrant question of \cite{FWI} at $\kappa$. The refutations known at
$\aleph_0$ \cite{AW,Koszmider} concern singular non-regular states and do not
transfer.
\end{question}

\begin{question}\label{q:continuous}
By \cite{BlecherWeaver}, Theorem 4.6, a $<\kappa$-additive pure state is normal on
every von Neumann subalgebra generated by fewer than $\kappa$ elements; in particular
any separably-acting chain such as $\{L^2[0,t]\}$ is annihilated and invisible to
$\FF_\varphi$. The $\kappa$-analogue of the continuous-masa question of \cite{FWII} is
therefore: working in $L^2(\{0,1\}^\kappa)$ with the product measure --- a
$\kappa$-generated diffuse masa, outside the reach of their Theorem 4.6 --- does the
family of spectral subspaces of the masa defeat every basis in the sense of
\cite{FWII}, Proposition 5.3? Theorem~\ref{thm:G} shows any witness family must be
transverse to every atomic masa; a diffuse masa is the extreme case.
\end{question}

\begin{question}\label{q:kappaC}
The $\kappa$-analogue of \cite{FWII}, Theorem 3.4, refined by
Proposition~\ref{prop:kappa-witness}: is $\Phi(\U)$ maximal if and only if $\U$ is a
$\sigma$-Q-point? Necessity is Corollary~\ref{cor:necessity}; sufficiency is open, and
would require an exact dichotomy (Theorem~\ref{thm:G}), or a substitute for it, proved
from the $\sigma$-Q-point property alone, without diagonal intersections.

Maximality does not, however, characterize normality up to isomorphism. In
\cite{FWIV} the Fubini product $\mathcal D\otimes\mathcal D$ of a normal measure with
itself is shown to satisfy exact paving outright --- Rowbottom homogeneity for finitely
many colorings of triples and quadruples replaces the diagonal intersection --- so that
$\Phi(\mathcal D\otimes\mathcal D)$ is a $\kappa$-complete nonprincipal pluf with
singleton face, while $\mathcal D\otimes\mathcal D$ is isomorphic to no normal
measure. The sharper questions are posed there: whether maximality, exact paving and
the $\sigma$-Q-point property coincide.
\end{question}

\begin{question}\label{q:calkinkappa}
The generalized Calkin algebra $\Blk/\mathcal K$ (or the quotient by operators of
support $<\kappa$) carries its own projection lattice. Do the maximal filters there
lift, in the sense of Farah--Wofsey's Question 6.39 \cite{FarahWofsey}, and does the
exact dichotomy of Theorem~\ref{thm:G} descend to the quotient along normal measures?
\end{question}

\subsection*{Acknowledgements}
This note redeems a suggestion made to the first author in a brief conversation, many
years ago, by the late Andrew Gleason: to consider these questions on a Hilbert space
whose dimension is a measurable cardinal. We are indebted to the paper of Blecher and
Weaver \cite{BlecherWeaver}, whose excision theory supplies the entire fourth section,
and to the survey of Farah and Wofsey \cite{FarahWofsey}. Portions of these results
were developed in the course of an extended dialogue with Anthropic's Claude language
model; responsibility for correctness rests with the authors.

\begin{thebibliography}{99}

\bibitem{AW} C.~Akemann and N.~Weaver, \emph{Not all pure states on $B(H)$ are
diagonalizable}, Proc. Natl. Acad. Sci. USA \textbf{105} (2008), 5313--5314.

\bibitem{Anderson79} J.~Anderson, \emph{Extensions, restrictions, and representations
of states on $C^*$-algebras}, Trans. Amer. Math. Soc. \textbf{249} (1979), 303--329.

\bibitem{BlecherWeaver} D.~P.~Blecher and N.~Weaver, \emph{Quantum measurable
cardinals}, arXiv:1607.08505.

\bibitem{FarahWofsey} I.~Farah and E.~Wofsey, \emph{Set theory and operator algebras},
in: Appalachian Set Theory 2006--2012 (J.~Cummings and E.~Schimmerling, eds.), London
Math. Soc. Lecture Note Ser. 406, Cambridge, 2013, 63--120.

\bibitem{FWI} D.~Feldman and A.~Wilce, \emph{Ultrafilter limits in the projection
lattice of a Hilbert space}, in preparation.

\bibitem{FWII} D.~Feldman and A.~Wilce, \emph{Intimate subspaces, block filters, and
the diagonalization of maximal projection filters}, in preparation.

\bibitem{FWIV} D.~Feldman and A.~Wilce, \emph{Exact paving for Fubini products of
normal measures, with an application to maximal projection filters}, in preparation.

\bibitem{Hamhalter} J.~Hamhalter, \emph{Quantum Measure Theory}, Fundamental Theories
of Physics 134, Kluwer, 2003.

\bibitem{KadisonSinger} R.~V.~Kadison and I.~M.~Singer, \emph{Extensions of pure
states}, Amer. J. Math. \textbf{81} (1959), 383--400.

\bibitem{Koszmider} P.~Koszmider, \emph{A non-diagonalizable pure state}, Proc. Natl.
Acad. Sci. USA \textbf{117} (2020), 33084--33089.

\bibitem{MSS} A.~Marcus, D.~A.~Spielman and N.~Srivastava, \emph{Interlacing families
II: Mixed characteristic polynomials and the Kadison--Singer problem}, Ann. of Math.
\textbf{182} (2015), 327--350.

\end{thebibliography}

\end{document}
